% =========================================================
% tag_appendix.tex
% Training-Adaptive Data Selection with a Reliability Gate (TAG)
% Appendices --- CIKM 2026 submission
%
% This appendix restores the full explanatory version of the
% trajectory-anchor appendix and retains the Reliability-Gate
% implementation used in the current TAG manuscript.
%
% Cross-references assumed to exist in the main paper:
%   \label{thm:anchor-stability}
%   \label{fig:anchor-drift}
%   \label{sec:anchor-stability}
%   \label{sec:experiments}
%
% Citation keys assumed (please verify against your .bib):
%   davis1970rotation, stewart1990matrix,
%   wang2023selfinstruct, xu2024wizardlm,
%   zhou2023lima, hendrycks2021mmlu, suzgun2022bbh,
%   cobbe2021gsm8k, patel2021svamp,
%   chen2021humaneval, austin2021mbpp,
%   clark2020tydiqa, artetxe2020xquad
% =========================================================


% ---------------------------------------------------------
\appendix


% =========================================================
% A. ANCHOR STABILITY
% =========================================================

\section{Proof of Theorem~\ref{thm:anchor-stability} and Empirical Envelope}
\label{app:anchor-stability-diagnostics}

This appendix proves Theorem~\ref{thm:anchor-stability}, clarifies the
cross-time eigengap convention, and reports the empirical
Davis--Kahan envelope used in Fig.~\ref{fig:anchor-drift}.

The result is deliberately local. It does not claim that the complete
data-selection ranking remains fixed over training. Instead, it shows
that when the representation covariance evolves smoothly and retains
a positive leading eigengap, the recomputed trajectory-anchor
direction changes at the optimizer scale. This is the stability
property needed for the anchor to act as a controlled modulation of
the dynamic carrier rather than as an unrelated source of variation.


\subsection{Proof of Theorem~\ref{thm:anchor-stability}}
\label{app:proof-stability}

\begin{proof}
Fix layer \(l\) and transition \(t-1\to t\). Write
\[
A=\Sigma_l^{(t-1)},\qquad
B=\Sigma_l^{(t)},\qquad
E=B-A .
\]
Let \(u=v_l^{(t-1)}\) be the top eigenvector of \(A\), and let
\(\tilde u\) be the uncalibrated top eigenvector of \(B\). Since
\[
B\tilde u=\lambda_1(B)\tilde u,
\]
the residual of \(\tilde u\) with respect to \(A\) is
\[
A\tilde u-\lambda_1(B)\tilde u
=
(A-B)\tilde u
=
-E\tilde u.
\]
Therefore,
\[
\|A\tilde u-\lambda_1(B)\tilde u\|_2
=
\|E\tilde u\|_2
\le
\|E\|_2
\le
\|E\|_F .
\]

By the residual form of the Davis--Kahan theorem
\cite{davis1970rotation,stewart1990matrix},
\[
\sin\Theta(u,\tilde u)
\le
\frac{\|E\|_F}
{\lambda_1(B)-\lambda_2(A)} .
\]
Using the assumptions of Theorem~\ref{thm:anchor-stability},
\[
\|\Sigma_l^{(t)}-\Sigma_l^{(t-1)}\|_F
\le
C_{\Sigma,l}\eta_{t-1}
\]
and
\[
\lambda_1(\Sigma_l^{(t)})
-
\lambda_2(\Sigma_l^{(t-1)})
\ge
\gamma_l,
\]
we obtain
\[
\sin\Theta
\le
\frac{C_{\Sigma,l}}{\gamma_l}\eta_{t-1}.
\]

After temporal sign calibration,
\[
\langle v_l^{(t)},v_l^{(t-1)}\rangle\ge0,
\]
so the corresponding principal angle satisfies
\[
\Theta\in[0,\pi/2].
\]
For unit vectors,
\[
\|v_l^{(t)}-v_l^{(t-1)}\|_2
=
2\sin(\Theta/2)
\le
2\sin\Theta .
\]
Therefore,
\[
\boxed{
\|v_l^{(t)}-v_l^{(t-1)}\|_2
\le
\frac{2C_{\Sigma,l}}{\gamma_l}\eta_{t-1}
}.
\]
\end{proof}


\paragraph{Eigengap convention.}
The theorem uses the cross-time gap
\[
\lambda_1(\Sigma_l^{(t)})
-
\lambda_2(\Sigma_l^{(t-1)})
\]
because the proof treats the new top eigenvector as an approximate
eigenvector of the previous covariance. By Weyl's inequality,
\[
\left|
\lambda_1(\Sigma_l^{(t)})
-
\lambda_1(\Sigma_l^{(t-1)})
\right|
\le
\|\Sigma_l^{(t)}-\Sigma_l^{(t-1)}\|_F .
\]
Under the covariance-drift assumption, the difference between the
cross-time gap and the usual same-time eigengap is therefore
\(O(\eta_{t-1})\).

The bound controls only the recomputed anchor direction, not the full
ranking or optimization trajectory. In TAG, this distinction is
especially useful because the reliability gate is static across
refreshes. The gate therefore introduces no additional temporal
variation into the anchor-stability argument.


\subsection{Empirical Davis--Kahan Tail}
\label{app:proof-envelope}

For Fig.~\ref{fig:anchor-drift}, we replace the worst-case constants
with realized per-transition quantities. Define
\[
\gamma_l^{(s)}
:=
\lambda_1(\Sigma_l^{(s+1)})
-
\lambda_2(\Sigma_l^{(s)}).
\]
For every transition satisfying
\[
\gamma_l^{(s)}>0,
\]
the one-step perturbation argument gives
\[
\|v_l^{(s+1)}-v_l^{(s)}\|_2
\le
\frac{
2\|\Sigma_l^{(s+1)}-\Sigma_l^{(s)}\|_F
}{
\gamma_l^{(s)}
}.
\]

Applying the triangle inequality from refresh \(t\) through a later
refresh \(T\) yields the cumulative envelope
\begin{equation}
\|v_l^{(t)}-v_l^{(T)}\|_2
\le
\sum_{s=t}^{T-1}
\frac{
2\|\Sigma_l^{(s+1)}-\Sigma_l^{(s)}\|_F
}{
\gamma_l^{(s)}
}.
\label{eq:tail-envelope}
\end{equation}

This is the blue dashed curve in Fig.~\ref{fig:anchor-drift}. It is an
upper envelope, not a fitted or tight predictor of measured anchor
drift. Its role is to compare the realized anchor movement with the
local perturbation scale implied by the covariance dynamics.


\subsection{Realized Stability Constants}
\label{app:stability-assumptions}

On the Qwen2.5-0.5B + LoRA run, computed over the \(1{,}496\) valid
post-warmup layer--refresh pairs, the realized covariance drift is
bounded by
\[
\sup_{l,t}
\frac{
\|\Sigma_l^{(t)}-\Sigma_l^{(t-1)}\|_F
}{
\eta_{t-1}
}
\le
C_\Sigma
=
5.78\times10^{3},
\]
and the retained pairs satisfy
\[
\inf_{l,t}
\left(
\lambda_1(\Sigma_l^{(t)})
-
\lambda_2(\Sigma_l^{(t-1)})
\right)
\ge
\gamma_{\min}
=
0.05 .
\]

The raw verifier eigenvectors required negative orientation correction
in \(1.7\%\) of the retained trajectory-anchor pairs. After sign
calibration, all reported displacements satisfy
\[
\langle v_l^{(t)},v_l^{(t-1)}\rangle\ge0 .
\]

With
\[
\eta=2\times10^{-4},
\]
the global plug-in envelope is
\[
(2C_\Sigma/\gamma_{\min})\eta
=
46.24 .
\]
This exceeds the maximum possible sign-calibrated step displacement
\[
\sqrt{2},
\]
so the global worst-case constant is numerically conservative for this
run. The informative comparison is therefore the realized envelope in
Eq.~\eqref{eq:tail-envelope}, which uses the observed covariance drift
and eigengap at each refresh.


% =========================================================
% B. HYPERPARAMETERS
% =========================================================

\section{Hyperparameter Configuration}
\label{app:hyperparams}

Table~\ref{tab:hyperparams} summarizes the per-backbone configuration
referenced from \S\ref{sec:experiments}. Within each backbone, all
selectors share these settings, so the selected subset is the only
varying factor across the methods we compare.

The reliability gate is computed before fine-tuning at the frozen base
checkpoint and cached for subsequent refreshes. Its implementation is
described separately in Appendix~\ref{app:reliability-gate}.

\begin{table}[h]
\centering
\caption{Hyperparameter settings used in our experiments.
Footnotes cover model-specific deviations.}
\label{tab:hyperparams}
\small
\setlength{\tabcolsep}{4pt}
\begin{tabular}{ll@{\hspace{1em}}ll}
\toprule
\multicolumn{2}{l}{\emph{Supervised fine-tuning}} &
\multicolumn{2}{l}{\emph{Dynamic Selection \& TAG}} \\
Parameter & Value & Parameter & Value \\
\midrule
Optimiser            & AdamW (8-bit)$^{a}$        & Selection ratio $\rho$  & 0.10 \\
Precision            & bf16                       & Anchor weight $\lambda$ & 1.0  \\
Learning rate        & $2\times10^{-5}$ $^{b}$   & Probe size $n_p$        & 1024 \\
LR scheduler         & cosine                     & Anchor layers           & all  \\
Warmup ratio         & 0.03 $^{c}$                & PCA batch size          & 4    \\
Weight decay         & 0.1                        & Episode batch size      & 16   \\
Gradient clip        & 1.0 $^{c}$                 & Gate checkpoint         & base \\
Gradient checkpoint  & enabled                    & Gate refresh            & none \\
Cutoff length        & 1024                       & Gate caching            & enabled \\
Epochs               & 3                          &                         &      \\
Effective batch size & 128 $^{d}$                 &                         &      \\
\bottomrule
\end{tabular}

\smallskip
\begin{flushleft}
\tiny
$^{a}$ Bitsandbytes 8-bit AdamW; auto-falls back to fp32 if the library is unavailable.\\
$^{b}$ Qwen2.5-7B uses $1\times10^{-5}$ per its official SFT recipe.\\
$^{c}$ LLaMA-2-7B + TAG uses warmup $0.06$ and gradient clip $0.5$ as a small SFT-side stability lever; all other (model, method) cells use the values shown.\\
$^{d}$ Per-GPU batch $8\times$ grad-accum $4\times$ DDP world-size $4$.
\end{flushleft}
\end{table}


\paragraph{Implementation note.}
The codebase contains optional actor--critic components retained for
future verifiable-reward extensions; they are not used in the ranking
rule reported in this paper. The dynamic component uses the
deterministic composite-reward carrier multiplied by the
trajectory-anchor factor. TAG prepends the static reliability gate,
giving the complete score
\[
s_i^{(t)}
=
G_i
\cdot
R_i^{(t)}
\left(
1+\lambda\,\mathrm{align}_i^{(t)}
\right).
\]
The gate \(G_i\) is computed once before fine-tuning, whereas
\(R_i^{(t)}\) and \(\mathrm{align}_i^{(t)}\) are refreshed as the
model moves through training.


% =========================================================
% C. RELIABILITY GATE
% =========================================================

\section{Reliability-Gate Implementation}
\label{app:reliability-gate}

The reliability gate addresses a different question from the two
training-adaptive utility views. Difficulty, uncertainty, and
trajectory alignment measure how useful a candidate appears at the
current model state. The gate instead asks whether the
instruction--response pair provides credible supervision in the first
place.

This distinction matters under low-quality instruction data. A
mismatched pair, a subtly wrong answer, or another corrupted example
may have high loss or uncertainty precisely because it is broken.
Without a separate reliability view, such examples can therefore be
preferred by a selector that interprets hardness only as
informativeness.

TAG separates the two decisions. Reliability is estimated once from
the frozen base model, and dynamic utility is subsequently evaluated
only through the training-adaptive carrier and trajectory anchor. The
complete fusion rule is
\begin{equation}
s_i^{(t)}
=
G_i
\cdot
R_i^{(t)}
\left(
1+\lambda\,\mathrm{align}_i^{(t)}
\right),
\label{eq:tag-score-app}
\end{equation}
where
\[
G_i\in[0,1]
\]
is static, while
\[
R_i^{(t)}
\quad\text{and}\quad
\mathrm{align}_i^{(t)}
\]
vary with the fine-tuning trajectory.


\subsection{Counterfactual Instruction Contrast}
\label{app:gate-counterfactual}

Consider instruction--response pair
\[
(x_i,y_i).
\]
Let
\[
\ell_k(y_i\mid x)
\]
denote the negative log-likelihood of the \(k\)-th response token
under the frozen base model when conditioned on instruction \(x\).

For every candidate \(i\), we construct a length-matched unrelated
instruction
\[
x_i^{-}
\]
sampled from the candidate pool. The response \(y_i\) is kept fixed.
The per-token counterfactual gain is
\begin{equation}
\delta_{i,k}
=
\ell_k(y_i\mid x_i^{-})
-
\ell_k(y_i\mid x_i).
\label{eq:gate-token-contrast}
\end{equation}

A positive value means that the observed instruction makes the
response easier to predict than the unrelated counterfactual
instruction. A value near zero means that replacing the instruction
has little effect on the response likelihood, providing little
evidence that the response is specifically supported by its observed
instruction.

The counterfactual is placed on the instruction side deliberately.
The objective is not to ask whether the response is fluent in
isolation. Instead, it asks whether the relationship between the
instruction and response is meaningful. This distinction is important
for mismatched pairs: a response can be perfectly fluent and
well-formed while answering a different instruction.


\subsection{Overall Relative Gain}
\label{app:gate-global-gain}

Summing the raw token contrasts alone would make the scale depend on
the intrinsic difficulty and length of the response. We therefore
normalize by the counterfactual loss.

Let
\[
L(y_i\mid x)
=
\sum_k \ell_k(y_i\mid x).
\]
The overall relative gain is
\begin{equation}
\bar{\Delta}_i
=
\frac{
\sum_k \delta_{i,k}
}{
\sum_k \ell_k(y_i\mid x_i^{-})
}
=
1-
\frac{
L(y_i\mid x_i)
}{
L(y_i\mid x_i^{-})
}.
\label{eq:gate-global-gain}
\end{equation}

This normalization makes the statistic scale-free. In particular, an
intrinsically difficult response that receives large loss under both
the true and counterfactual instructions is not automatically assigned
a large reliability gain.


\subsection{Span-Level Tail Gain}
\label{app:gate-span-gain}

An average response-level statistic can hide localized corruption.
For example, a response may be largely correct but contain an
incorrect final answer or a locally spliced segment. Most tokens can
still benefit from the true instruction, causing the response-level
mean to remain positive even though a critical region is unreliable.

To expose such localized failures, the response is partitioned into
contiguous spans
\[
\mathcal S_1,\ldots,\mathcal S_{M_i},
\]
implemented as sentences or fixed windows of \(W\) tokens. For span
\(\mathcal S_m\), define
\begin{equation}
\Delta_{i,m}
=
1-
\frac{
\sum_{k\in\mathcal S_m}
\ell_k(y_i\mid x_i)
}{
\sum_{k\in\mathcal S_m}
\ell_k(y_i\mid x_i^{-})
}.
\label{eq:gate-span-gain}
\end{equation}

Not every span contains enough instruction-dependent content to be
meaningful. Let
\[
\mathcal C_i
\]
denote the set of eligible spans. A span is excluded when its
counterfactual loss
\[
\sum_{k\in\mathcal S_m}
\ell_k(y_i\mid x_i^{-})
\]
falls below the frozen low-information threshold. This prevents
boilerplate such as short acknowledgements or closing pleasantries
from controlling the reliability decision merely because such text
is naturally weakly dependent on the instruction.

The tail statistic is the least-supported eligible span:
\begin{equation}
\Delta_i^{\min}
=
\min_{m:\mathcal S_m\in\mathcal C_i}
\Delta_{i,m}.
\label{eq:gate-tail-gain}
\end{equation}

The gate therefore requires evidence at two resolutions: the
instruction should help the response overall, and it should also help
the least-supported meaningful span.


\subsection{Completeness Check and Zero-Anchored Gate}
\label{app:gate-final}

The likelihood contrast does not fully capture truncation. If a
response is cut off prematurely, every token that remains may still
be well supported by the true instruction. TAG therefore includes a
raw-text completeness term
\[
c_i\in\{1,c_{\mathrm{trunc}}\},
\]
where
\[
c_i=1
\]
for a response that terminates properly and
\[
c_i=c_{\mathrm{trunc}}
\]
otherwise.

The two likelihood views are combined conservatively:
\begin{equation}
\widehat{\Delta}_i
=
\min
\left(
\bar{\Delta}_i,
\Delta_i^{\min}
\right).
\label{eq:gate-combined-gain}
\end{equation}

Let \(\sigma(\cdot)\) denote the logistic function and let \(s>0\)
denote the gate scale calibrated once per backbone on the clean
reference pool. The final reliability gate is
\begin{equation}
G_i
=
c_i
\left(
2\sigma
\left(
\frac{\widehat{\Delta}_i}{s}
\right)
-1
\right)_{+},
\label{eq:reliability-gate}
\end{equation}
where
\[
(a)_{+}:=\max(a,0).
\]

Equation~\eqref{eq:reliability-gate} is zero-anchored. Since
\[
\sigma(0)=\frac{1}{2},
\]
we have
\[
\widehat{\Delta}_i\le0
\quad\Longrightarrow\quad
G_i=0.
\]

Thus, ``no gain'' is an absolute boundary rather than an arbitrary
pool-relative rank. An example whose observed instruction provides
no positive evidence over its unrelated counterfactual receives an
exactly zero reliability value.


\subsection{Why Reliability Is a Veto Rather Than a Vote}
\label{app:gate-veto}

The multiplicative fusion in Eq.~\eqref{eq:tag-score-app} is
non-compensatory. If
\[
G_i=0,
\]
then
\[
s_i^{(t)}=0
\qquad
\text{for every refresh }t.
\]

Consequently, high difficulty, uncertainty, or trajectory alignment
cannot rescue a candidate that fails the reliability gate. This is
intentional. Dynamic utility can decide which reliable examples are
most useful at the current training state, but it cannot manufacture
credible supervision from an unreliable instruction--response pair.

For candidates satisfying
\[
G_i>0,
\]
selection remains fully training-adaptive through
\[
R_i^{(t)}
\left(
1+\lambda\,\mathrm{align}_i^{(t)}
\right).
\]
The gate therefore does not replace the dynamic selector. It defines
how strongly each candidate is allowed to participate in that
selector.


\subsection{Why the Counterfactual Changes the Instruction}
\label{app:gate-why-instruction}

At first sight, perturbing the instruction may appear indirect because
many forms of corruption are introduced on the response side. The
choice is deliberate.

First, mismatched supervision breaks the dependency between the
instruction and the response rather than necessarily damaging the
surface form of the response. A mismatched response can remain fluent,
grammatical, and semantically coherent. A response-only fluency or
perplexity criterion can therefore treat it as clean.

By holding \(y_i\) fixed and replacing \(x_i\) with \(x_i^{-}\), the
gate directly asks whether the model's probability of that response
depends on the observed instruction. If the likelihood remains almost
unchanged after replacing the instruction, the pair contains little
instruction-specific dependency worth training on.

Second, a response-side contrast tends to collapse back into a
fluency criterion. Comparing a response with a damaged version of
itself primarily tests whether the original response is more typical
or predictable than the perturbation. That information substantially
overlaps with loss and perplexity, while remaining weak against fluent
mismatches and fluent but incorrect answers.

Third, the instruction side supplies a counterfactual without an
external oracle. Constructing a ``correct'' alternative response would
require knowing or generating the correct answer, introducing an
external judge, stronger model, or other supervision. In contrast,
\(x_i^{-}\) can be sampled directly from the existing pool.

The injected corruption used for evaluation should not be confused
with the detection mechanism. The experimenter may know how an
example was corrupted in order to measure detection accuracy, but the
gate itself operates without access to that corruption label or
mechanism.


\subsection{Base-Checkpoint Computation and Caching}
\label{app:gate-cache}

The reliability gate is a property of the data pair rather than of
the evolving training trajectory. It is therefore evaluated at the
frozen base checkpoint before fine-tuning begins.

For each example, the implementation performs the following steps:
\begin{enumerate}
    \item obtain the observed instruction \(x_i\) and response \(y_i\);
    \item sample a length-matched unrelated instruction \(x_i^{-}\);
    \item evaluate token losses under \(x_i\) and \(x_i^{-}\);
    \item compute the token contrasts in
          Eq.~\eqref{eq:gate-token-contrast};
    \item compute the overall relative gain in
          Eq.~\eqref{eq:gate-global-gain};
    \item construct eligible response spans and compute
          Eq.~\eqref{eq:gate-tail-gain};
    \item apply the raw-text completeness check \(c_i\);
    \item compute the zero-anchored gate in
          Eq.~\eqref{eq:reliability-gate}; and
    \item cache \(G_i\) with candidate \(i\).
\end{enumerate}

After this initial profiling pass, \(G_i\) is not recomputed during
fine-tuning. Each dynamic refresh retrieves the cached gate and
recomputes only the training-state-dependent carrier and
trajectory-anchor terms.

Hence, the reliability view adds an initialization-time screening
operation but no per-refresh model evaluation. It also adds no new
temporal instability to the trajectory-anchor analysis because
\(G_i\) remains constant over \(t\).


% =========================================================
% D. RANDOM SELECTION-RATIO SWEEP
% =========================================================

\begin{table}[t]
\centering
\caption{\textbf{Random selection-ratio sweep on LLaMA-2-7B +
Alpaca-GPT4.} Budget-only controls with all other hyperparameters
held fixed; W-AVG plateaus within a narrow range across
\(\rho\in\{10,20,30,40,50\}\%\).}
\label{tab:ratio-sweep}
{\fontsize{7.5pt}{9pt}\selectfont
\setlength{\tabcolsep}{1.8pt}
\renewcommand{\arraystretch}{1.05}
\begin{tabular}{@{}lcccccccc c@{}}
\toprule
& \multicolumn{2}{c}{Know.} & \multicolumn{2}{c}{Math}
& \multicolumn{2}{c}{Code} & \multicolumn{2}{c}{MLing.} & \\
\cmidrule(lr){2-3}
\cmidrule(lr){4-5}
\cmidrule(lr){6-7}
\cmidrule(lr){8-9}
$\rho$ & MMLU & BBH & SVAMP & GSM8K & MBPP & H-Eval & TydiQA & XQuAD & W-AVG \\
\midrule
10\% & 45.04 & 35.54 & 36.80 & 13.57 & 35.41 & 24.51 & 60.83 & 52.28 & 38.00 \\
20\% & 45.53 & 35.72 & 38.33 & 15.92 & 34.24 & 24.13 & 61.31 & 53.31 & 38.56 \\
30\% & 45.91 & 36.23 & 37.33 & 14.10 & 36.96 & 25.31 & 61.32 & 53.73 & 38.86 \\
40\% & 46.21 & 36.37 & 40.00 & 14.33 & 35.02 & 24.02 & 61.29 & 53.37 & 38.83 \\
50\% & 46.31 & 36.42 & 38.00 & 13.95 & 33.85 & 25.64 & 60.98 & 53.67 & 38.60 \\
\bottomrule
\end{tabular}
}
\end{table}


\section{Random Selection-Ratio Sweep}
\label{app:ratio-sweep}

The random rows in Table~\ref{tab:ratio-sweep} are
\emph{budget-only controls}: all hyperparameters except \(\rho\) are
fixed, so differences across rows reflect data budget rather than
selection quality.

W-AVG changes only within a \(0.86\)-pp band from \(\rho=10\%\) to
\(\rho=50\%\), peaking at \(\rho=30\%\) before saturating. Thus,
increasing the random budget up to \(5\times\) provides little
additional headroom.

The task-level behavior supports the same interpretation. Performance
does not increase monotonically with the number of randomly selected
examples. Some mathematical tasks improve at intermediate ratios,
whereas several code and multilingual scores fluctuate around a
similar range. The principal effect is therefore saturation rather
than a consistent benefit from additional random data.

We use \(\rho=10\%\) throughout the main comparison to test whether a
selection signal can improve training efficiency under a constrained
budget. If TAG at \(\rho=10\%\) exceeds random selection even when the
random baseline uses substantially more data, the difference cannot be
explained simply by the number of selected examples.

Each ratio uses one fixed-seed random draw, so across-seed variance is
not estimated by this diagnostic. The purpose of the sweep is narrower:
to test whether the operating budget lies in a regime where simply
increasing random data provides substantial additional headroom. The
observed saturation is consistent with prior instruction-tuning
findings that data quality can matter more than raw scale once a modest
subset is used~\cite{zhou2023lima}.


% =========================================================
% E. DATASETS AND BENCHMARK SCORING
% =========================================================

\section{Datasets and Benchmark Scoring}
\label{app:benchmarks}

This appendix describes the instruction-tuning data and benchmark
scoring protocols. All benchmarks are evaluated with
\texttt{lm\allowbreak-evaluation\allowbreak-harness}. Unless otherwise
stated, AVG denotes the unweighted mean over the evaluated tasks.

The downstream suite covers knowledge and general reasoning,
mathematical reasoning, code generation, and multilingual question
answering. The same evaluation configuration is used for every
selection method within a backbone.


\subsection{Instruction-Tuning Pools}
\label{app:instruction-pools}

\paragraph{Alpaca-GPT4.}
Alpaca-GPT4 contains approximately 52K Stanford-Alpaca-style
instructions paired with GPT-4 responses, covering broad open-domain
instruction following~\cite{wang2023selfinstruct}. The pool includes
factual questions, explanation, rewriting, reasoning, and lightweight
programming tasks.

Its breadth makes it useful for evaluating a selector whose goal is
not to optimize one downstream benchmark but to identify broadly
useful supervision from a heterogeneous instruction pool.


\paragraph{WizardLM Evol-Instruct.}
WizardLM Evol-Instruct contains evolution-generated instructions
designed to be harder, more diverse, and more reasoning-heavy
\cite{xu2024wizardlm}. Its construction progressively modifies
instructions to increase complexity, reasoning depth, specificity, or
constraint structure.

We use this pool as a complementary instruction distribution rather
than as a source of method-specific tuning. The conceptual roles of
the reliability gate, dynamic carrier, and trajectory anchor remain
unchanged.


\subsection{Knowledge and General Reasoning}
\label{app:benchmark-knowledge}

\paragraph{MMLU.}
MMLU is a multiple-choice benchmark spanning academic and professional
subjects, including science, mathematics, humanities, social science,
law, medicine, and other domains~\cite{hendrycks2021mmlu}. Its score
is answer-choice accuracy.

MMLU provides a broad test of whether instruction tuning preserves or
improves general-purpose knowledge use rather than specializing only
to the selected training examples.


\paragraph{BBH.}
BIG-Bench Hard (BBH) is a suite of challenging symbolic, logical,
algorithmic, and commonsense reasoning tasks
\cite{suzgun2022bbh}. Its score is exact-match accuracy under the
common evaluation harness.

BBH complements MMLU by placing greater emphasis on explicit
multi-step reasoning rather than broad disciplinary knowledge.


\subsection{Mathematical Reasoning}
\label{app:benchmark-math}

\paragraph{GSM8K.}
GSM8K contains grade-school mathematical word problems requiring
multi-step arithmetic reasoning~\cite{cobbe2021gsm8k}. The model
generates a solution and final numeric answer. We normalize the final
answer and report exact-match accuracy.


\paragraph{SVAMP.}
SVAMP evaluates arithmetic word-problem reasoning under variations in
surface form and problem structure~\cite{patel2021svamp}. It is
designed to test whether models capture the underlying quantitative
relation rather than relying on shallow templates.

As with GSM8K, the final normalized numeric answer is scored by exact
match.


\subsection{Code Generation}
\label{app:benchmark-code}

\paragraph{HumanEval.}
HumanEval contains Python programming problems specified by a natural
language description, function signature, and hidden unit tests
\cite{chen2021humaneval}. The model must generate executable code that
satisfies the tests.

We report pass@1 under the fixed execution-based evaluation protocol.


\paragraph{MBPP.}
Mostly Basic Python Problems (MBPP) evaluates short Python programs
from natural-language task descriptions and examples
\cite{austin2021mbpp}. Evaluation is execution-based rather than
based on lexical similarity to a reference program.

We report pass@1. Using both MBPP and HumanEval reduces dependence on
the characteristics of a single code-generation benchmark.


\subsection{Multilingual Question Answering}
\label{app:benchmark-multilingual}

\paragraph{TyDiQA.}
TyDiQA evaluates question answering across typologically diverse
languages and was constructed to reduce the English-centric bias of
many multilingual QA resources~\cite{clark2020tydiqa}. Each example
contains a question and a supporting context passage.

We report normalized token-overlap F1 averaged over the evaluated
languages.


\paragraph{XQuAD.}
XQuAD provides parallel SQuAD-style question-answering examples across
multiple languages~\cite{artetxe2020xquad}. It complements TyDiQA with
a controlled cross-lingual benchmark based on aligned question and
context examples.

We again report normalized token-overlap F1 averaged over languages.


\subsection{Aggregate Reporting}
\label{app:benchmark-aggregation}

For descriptive reporting, the eight benchmarks form four capability
groups:
\[
\begin{aligned}
\text{Knowledge} &: \{\text{MMLU},\text{BBH}\},\\
\text{Math} &: \{\text{SVAMP},\text{GSM8K}\},\\
\text{Code} &: \{\text{MBPP},\text{HumanEval}\},\\
\text{Multilingual} &: \{\text{TyDiQA},\text{XQuAD}\}.
\end{aligned}
\]

The task roster and scoring rules are fixed across selectors. No
benchmark is added, removed, or reweighted based on the performance of
a particular selection method.

The reliability gate and dynamic selection score are computed entirely
from the instruction-tuning pool and the model state. Downstream
benchmark labels and outcomes are never used to construct the
selection ranking.


% =========================================================
% END APPENDIX
% =========================================================
